A finals de l'any 1933, a la Universitat de Stanford (EUA), Fritz Zwicky i Walter Baade van proposar per primera vegada l'existència de les estrelles de neutrons. Aquestes estrelles, formades només per neutrons, es poden originar després de l'explosió d'una supernova. Els neutrons que les formen són el resultat de la fusió de protons i electrons, provocada per la compressió que exerceix el camp gravitatori d'aquestes estrelles. Per a una estrella de neutrons determinada que té una massa de $2,9\times 10^{30}\ \mathrm{kg}$ i un radi de 10~km, calculeu:

\figura[0.4]{fig1.png}

\begin{apartat}{1}
El mòdul de la intensitat de camp gravitatori que l'estrella de neutrons crea a la seva pròpia superfície.
\end{apartat}

\begin{apartat}{1}
La velocitat mínima que hem de donar a un coet en el moment del llançament des de la superfície de l'estrella perquè es pugui escapar de l'atracció d'aquesta (ignoreu els possibles efectes relativistes). Demostreu l'expressió utilitzada per a fer el càlcul i feu esment del principi de conservació en què us baseu.
\end{apartat}

\blocetiquetat{Dada:}{$G = 6,67\times 10^{-11}\ \mathrm{N\,m^2\,kg^{-2}}$.}
